%%%%%%%%%%%%%%%%%%%%%%%%%%%%%%%%%%%%%%%%%%%%%%%%%%%%%%%%%%%
% CheckPostconditionSatisfiability — VerCors Report
%%%%%%%%%%%%%%%%%%%%%%%%%%%%%%%%%%%%%%%%%%%%%%%%%%%%%%%%%%%

\documentclass[9pt,a4paper,twocolumn]{article}
\usepackage[english]{babel}
\usepackage[T1]{fontenc}
\usepackage[utf8]{inputenc}
\usepackage{listings}
\usepackage{xcolor}
\usepackage{booktabs}
\usepackage{array}
\usepackage{geometry}
\usepackage{titlesec}
\usepackage{parskip}
\usepackage{microtype}
\usepackage{float}
\usepackage[hidelinks]{hyperref}

\geometry{a4paper, margin=1.8cm, columnsep=0.6cm}

\lstset{
  basicstyle=\ttfamily\small,
  frame=single,
  breaklines=true,
  columns=fullflexible,
  keepspaces=true,
  xleftmargin=4pt,
  xrightmargin=4pt
}

\titleformat{\section}{\normalfont\large\bfseries}{\thesection}{0.8em}{}
\titleformat{\subsection}{\normalfont\normalsize\bfseries}{\thesubsection}{0.8em}{}
\titleformat{\subsubsection}{\normalfont\normalsize\itshape}{\thesubsubsection}{0.8em}{}
\titlespacing{\section}{0pt}{8pt}{4pt}
\titlespacing{\subsection}{0pt}{6pt}{2pt}
\titlespacing{\subsubsection}{0pt}{4pt}{2pt}

\setlength{\parskip}{4pt}
\setlength{\parindent}{0pt}

\title{\textbf{Postcondition Satisfiability Checking in VerCors}}
\author{Alexandra Jakovleva}
\date{\today}

\begin{document}

\maketitle

%----------------------------------------------------------
\section{Introduction}

A postcondition that is internally contradictory — for example
\texttt{ensures x > 0 \&\& x < 0} — can never hold in any
program state. The verifier will prove the method correct, because
from a false postcondition anything is derivable, but the contract
is meaningless: no execution of the method can ever satisfy it.
This is a specification bug that currently produces no warning.

\texttt{CheckPostconditionSatisfiability} addresses this by
generating an isolated check for each method, function, and
procedure that proves the postcondition is satisfiable — that
some state exists in which it can hold. If no such state exists,
a \texttt{NontrivialPostconditionUnsatisfiable} error is reported
directly on the contract.

\textbf{Scope.} The pass checks the \texttt{ensures} clause of
every \texttt{ContractApplicable} node: methods, procedures, and
abstract functions. It does not check preconditions; those are
covered by the existing \texttt{CheckContractSatisfiability} pass.
Trivially true postconditions (\texttt{true}) are skipped as
uninteresting. Trivially false postconditions (\texttt{false})
are also skipped because the verifier already rejects them
unconditionally; no additional check is needed.

%----------------------------------------------------------
\section{What It Solves}

The most direct case is a self-contradictory boolean postcondition:

\begin{lstlisting}
ensures x > 0 && x < 0; // unsatisfiable
void f(int x) { ... }
\end{lstlisting}

Without the check, VerCors verifies \texttt{f} successfully and
silently. With it, \texttt{NontrivialPostconditionUnsatisfiable}
fires on the contract.

A less obvious case arises with resource postconditions. A
postcondition such as \texttt{ensures Perm(x, write) **
Perm(x, write)} cannot hold in any heap state — write permission
on the same location cannot appear twice. The satisfiability
check catches this in isolation, before verification of the method
body begins.

%----------------------------------------------------------
\section{Key Design Decision: Postcondition in Isolation}

A natural first attempt would be to check the postcondition
together with the precondition: if \texttt{requires P; ensures Q},
check whether \texttt{P \&\& Q} is satisfiable. This would produce
false positives. Consider:

\begin{lstlisting}
requires x == 5;
ensures x < 0;
void f(int x) { x = -1; }
\end{lstlisting}

Here \texttt{x == 5 \&\& x < 0} is unsatisfiable, yet the method
is perfectly correct — the body is allowed to change \texttt{x}.
The precondition describes the \emph{pre-state} and the
postcondition describes the \emph{post-state}; combining them as
if they applied to the same state is unsound.

The pass therefore checks the postcondition alone: it asks only
whether some post-state exists in which the postcondition holds,
with no reference to the precondition. This avoids the false
positive above while still catching genuinely unsatisfiable
postconditions.

%----------------------------------------------------------
\section{Implementation}

For each contract, the pass generates a standalone helper
procedure:

\begin{lstlisting}[language=Scala]
procedure(
  requires = UnitAccountedPredicate(postPred),
  body     = Some(Assert(ff)),
  ...
)
\end{lstlisting}

The postcondition \texttt{postPred} is placed in the
\texttt{requires} position of the helper. An \texttt{assert false}
forms the entire body, wrapped in an \texttt{ExpectedError} that
expects this assert to fail.

The reasoning is: if the postcondition is satisfiable, some
pre-state of the helper satisfies it, the helper is entered, and
\texttt{assert false} fires — the expected error is tripped and
verification passes. If the postcondition is unsatisfiable, no
pre-state can satisfy the \texttt{requires}, the helper is
unreachable, \texttt{assert false} never fires, the expected
error is not tripped, and \texttt{NontrivialPostconditionUnsatisfiable}
is reported.

Occurrences of \texttt{\textbackslash result} in the postcondition
are replaced by a fresh variable of the appropriate return type,
so the expression is well-formed in \texttt{requires} position.
The helper procedure is extracted via \texttt{Extract} and
declared as a top-level global procedure, fully isolated from
the original method's context.

The pass runs late in the Transformation pipeline, immediately
after \texttt{CheckContractSatisfiability}, once all
language-specific constructs have been lowered.

%----------------------------------------------------------
\section{Limitations}

\textbf{Postcondition only.} The check does not consider whether
the postcondition is \emph{achievable} given the precondition and
the method body. A postcondition that is satisfiable in isolation
but unreachable given the precondition is not caught. Detecting
this would require reasoning about the method body, which is
beyond the scope of a static syntactic check.

\textbf{Quantified postconditions.} If the postcondition contains
a quantifier whose body is trivially false, the check may not fire
depending on how the SMT solver handles the formula. The check
is sound — it never produces a false positive — but may miss some
unsatisfiable postconditions expressed through complex quantified
formulas.

\textbf{Well-formedness.} If the postcondition is not
well-formed (e.g.\ it contains a permission expression that is
syntactically malformed), the well-formedness error is silently
suppressed rather than blocking the satisfiability check.
\texttt{NotWellFormedIgnoreCheckPostSat} redirects such errors
to the \texttt{ExpectedError} handler so they do not produce
spurious failures.

%----------------------------------------------------------
\section{Testing Plan}

\begin{itemize}
  \item \textbf{Contradictory boolean postcondition.}
    \texttt{ensures x > 0 \&\& x < 0} — must fire
    \texttt{NontrivialPostconditionUnsatisfiable}.

  \item \textbf{Trivially false postcondition.}
    \texttt{ensures false} — must be skipped silently
    (already an unconditional verifier error).

  \item \textbf{Trivially true postcondition.}
    \texttt{ensures true} — must be skipped silently.

  \item \textbf{Satisfiable postcondition.}
    \texttt{ensures x > 0} with no precondition — must
    pass without warning.

  \item \textbf{Postcondition with \texttt{\textbackslash result}.}
    \texttt{ensures \textbackslash result > 0} on a
    non-void method — must handle \texttt{\textbackslash result}
    substitution and pass without warning.

  \item \textbf{False positive guard.}
    \texttt{requires x == 5; ensures x < 0} on a method
    that assigns \texttt{x = -1} — must \emph{not} fire,
    confirming isolation from the precondition.

  \item \textbf{Unsatisfiable resource postcondition.}
    \texttt{ensures Perm(x, write) ** Perm(x, write)} —
    must fire.

  \item \textbf{Void method.}
    A void method with a contradictory postcondition —
    must fire without attempting \texttt{\textbackslash result}
    substitution.
\end{itemize}

\end{document}
